\section{System Architecture} \label{sec:system-architecture}
The system is divided into two main components: a Python backend and a Vue frontend.
The backend is responsible for the formal SAT encoding, solver execution, model decoding, and output generation.
The frontend is responsible for user interaction, input validation, and visualization of the computed plan.

\begin{figure}[H]
    \centering
    \fbox{
        \begin{minipage}{0.9\textwidth}
            \centering
            User Input (Frontend or CLI)\\
            $\downarrow$\\
            Problem Configuration\\
            $\downarrow$\\
            SAT Encoder\\
            $\downarrow$\\
            CNF Formula + Variable Mapping\\
            $\downarrow$\\
            Glucose3 SAT Solver\\
            $\downarrow$\\
            Model Decoder\\
            $\downarrow$\\
            Human-Readable Plan + Visualization
        \end{minipage}
    }
    \caption{General pipeline of the Movers SAT Solver.}
    \label{fig:pipeline}
\end{figure}

\subsection{Backend Architecture} \label{subsec:backend-architecture}

The backend is implemented in Python using FastAPI\@.
It exposes REST endpoints that allow the frontend to send a problem instance and receive either a satisfiable plan or an unsatisfiable result.

\noindent The backend is organized into the following main modules:

\begin{itemize}
    \item \texttt{models/problem.py}: defines the Pydantic models for problem inputs and solver outputs.
    \item \texttt{encoding/variables.py}: defines the Boolean variable families and maps symbolic variables to integer identifiers.
    \item \texttt{encoding/common\_clauses.py}: provides reusable CNF helpers such as at-most-one, at-least-one, and exactly-one constraints.
    \item \texttt{encoding/task1\_encoder.py}: encodes the movement-only problem.
    \item \texttt{encoding/task2\_encoder.py}: extends Task 1 with item movement and carrying constraints.
    \item \texttt{solver/runner.py}: runs Glucose3 either for a fixed time horizon or with iterative deepening.
    \item \texttt{solver/decoder.py}: converts a satisfying SAT assignment into a readable sequence of time steps.
    \item \texttt{output/dimacs.py}: serializes the generated CNF formula in DIMACS format.
    \item \texttt{output/mapping.py}: exports the variable mapping as JSON\@.
    \item \texttt{api/routes.py}: defines the REST API endpoints.
\end{itemize}

\noindent The API provides endpoints for solving a problem optimally, solving at a fixed horizon, downloading the DIMACS CNF file, and downloading the variable mapping.

\subsection{Frontend Architecture} \label{subsec:frontend-architecture}

The frontend is implemented using Vue and Tailwind CSS\@.
It communicates with the backend through a small API client that sends JSON payloads to the FastAPI server and retrieves solver results.

\noindent The frontend is divided into several components:

\begin{itemize}
    \item \texttt{App.vue}: maintains the global form state and coordinates solver requests.
    \item \texttt{ProblemSetup.vue}: provides the form where the user selects the task type, number of floors, number of movers, items, and advanced solver options.
    \item \texttt{ResultsPanel.vue}: displays solver status, runtime, number of steps, playback controls, and download buttons.
    \item \texttt{BuildingVisualizer.vue}: renders the building floors and places movers and items on the correct floors for the selected time step.
    \item \texttt{ActionLog.vue}: lists the actions performed by each mover at the current time step.
\end{itemize}

\noindent The interface allows two solving modes:

\begin{itemize}
    \item \textbf{Iterative Deepening:} searches for the shortest valid solution up to a maximum number of steps.
    \item \textbf{Fixed Horizon:} checks whether a solution exists for an exact number of time steps.
\end{itemize}

\noindent When a solution is found, the user can play through the sequence using timeline controls.
The frontend allows the user to download the generated DIMACS CNF instance and the variable mapping.
The satisfying model returned by the SAT solver is decoded into the displayed step-by-step plan.
